\section{Noetherian Modules and PID Structure Theorems}

This subsection begins with the structure facts that make free modules and
finite generation usable: rank, splitting maps onto free modules, Noetherian
modules. It then develops the classification of
finitely generated modules over a principal ideal domain. The guiding model is
the abelian-group theorem over $\Z$ and, later, the polynomial-module
classification behind rational canonical form.

\subsection{Free Modules and Rank}

\begin{definition} \label{def:free-module-basis} \hlt{Free Module and Basis}
\begin{enumerate}[label=(\roman*)]
\setlength{\itemsep}{0pt}
\item A subset $\mathcal{B}\subseteq M$ is a \hlt{basis} of $M$ if $\mathcal{B}$ is $R$-linearly independent and $\langle\!\langle\mathcal{B}\rangle\!\rangle_R=M$.
\item An $R$-module $M$ is \hlt{free} if $M=\{0\}$ or $M$ has a basis.
\item If $\mathcal{B}$ is a basis of $M$, then $M$ is \hlt{free on $\mathcal{B}$}.
\end{enumerate}
\end{definition}

\begin{theorem} \label{thm:module-basis-unique-coordinates} \hlt{Basis Means Unique Finite Coordinates}\\
A subset $\mathcal{B}\subseteq M$ is a basis if and only if $\forall v\in M,\ \exists!\,(F,(r_b)_{b\in F})$, with $F\subseteq_{\operatorname{fin}}\mathcal{B},\ r_b\in R\setminus\{0\}$, $v=\sum\limits_{b\in F}r_bb$.\\
The case $F=\emptyset$ gives the empty sum $0$.
\begin{proof}
\begin{enumerate}[label=\roman*.]
\setlength{\itemsep}{0pt}
\item[$(\Rightarrow)$] Spanning gives a finite expression; collecting equal basis elements and deleting zero coefficients gives the displayed form. If $\sum\limits_{b\in F}r_bb=\sum\limits_{b\in G}s_bb$, then $\sum\limits_{b\in F\cup G}(r_b-s_b)b=0$, where missing coefficients are $0$.\\
Independence gives $r_b=s_b$ for all $b$, hence $F=G$ and the coefficients agree.
\item[$(\Leftarrow)$] Existence gives $\langle\!\langle\mathcal{B}\rangle\!\rangle_R=M$. If $\sum\limits_{b\in F}r_bb=0$ and $F_0\defeq\{b\in F:r_b\neq0\}$ is nonempty, then $0=\sum\limits_{b\in F_0}r_bb$ contradicts uniqueness of the empty expression. Thus all $r_b=0$, so $\mathcal{B}$ is independent.
\end{enumerate}
\hlt{Structural Commentary}: Bases give unique finite coordinates.\\
\hlt{Key Insight}: Existence is spanning; uniqueness is independence.
\end{proof}
\end{theorem}

\begin{proposition} \label{prop:module-basis-minimal-maximal} \hlt{Bases Are Minimal Spanning and Maximal Independent}\\
If $\mathcal{B}$ is a basis of $M$, then $\forall b\in\mathcal{B}$, $\langle\!\langle\mathcal{B}\setminus\{b\}\rangle\!\rangle_R\neq M$, and $\forall v\in M\setminus\mathcal{B}$, $\mathcal{B}\cup\{v\}$ is $R$-linearly dependent.
\begin{proof}
If $b\in\mathcal{B}$ and $b\in\langle\!\langle\mathcal{B}\setminus\{b\}\rangle\!\rangle_R$, then $b$ has two finite coordinate expressions, contradicting Theorem \ref{thm:module-basis-unique-coordinates}.\\
Hence deleting $b$ destroys spanning. If $v\notin\mathcal{B}$, write $v=\sum_{b\in F}r_bb$ with $F\subseteq_{\operatorname{fin}}\mathcal{B}$.\\
Then $v-\sum\limits_{b\in F}r_bb=0$ is a nontrivial relation in $\mathcal{B}\cup\{v\}$, so $\mathcal{B}\cup\{v\}$ is dependent.\\
\hlt{Structural Commentary}: Basis still implies the familiar minimal/maximal properties.\\
\hlt{Key Insight}: The converses are vector-space privileges, not general module facts.
\end{proof}
\end{proposition}

\begin{example} \label{ex:free-module-with-nonfree-submodule} \hlt{A Free Module with a Nonfree Submodule}\\
Let $R=\mathbb{Z}\times\mathbb{Z}$ with componentwise multiplication. As an $R$-module, $R=R(1,1)$, so $R$ is free on $\{(1,1)\}$.\\
Set $N\defeq\mathbb{Z}\times\{0\}\leq R$. For every $0\neq(a,0)\in N$, $(0,1)(a,0)=0$.\\
Thus every nonzero element of $N$ is torsion, so $N$ has no nonempty independent subset and is not free.
\end{example}

\begin{perspective} \label{per:dimension-as-field-case-of-rank} \hlt{Dimension as the Field Case of Rank}\\
An $\mathbb{F}$-vector space is a free $\mathbb{F}$-module. By Theorem \ref{thm:basis-length-well-defined}, finite bases have invariant cardinality; over general rings, rank invariance requires additional hypotheses.
\end{perspective}

\begin{theorem} \label{thm:free-module-map-extension} \hlt{Maps Out of a Free Module}\\
Let $M$ be free on $\mathcal{B}$ and let $N$ be an $R$-module.\\
For every function $f:\mathcal{B}\to N$, there exists a unique $\widetilde{f}\in\operatorname{Hom}_R(M,N)$ such that $\widetilde{f}|_{\mathcal{B}}=f$.
\begin{proof} For $v=\sum\limits_{b\in F}r_bb$ with $F\subseteq_{\operatorname{fin}}\mathcal{B}$, define $\widetilde{f}(v)\defeq\sum\limits_{b\in F}r_bf(b)$.\\
Theorem \ref{thm:module-basis-unique-coordinates} makes the expression of $v$ unique, so $\widetilde{f}$ is well-defined.\\
For $x=\sum a_bb$ and $y=\sum c_bb$, combining supports gives $\widetilde{f}(rx+sy)=r\widetilde{f}(x)+s\widetilde{f}(y)$, so $\widetilde{f}$ is $R$-linear and clearly extends $f$. If $\phi\in\operatorname{Hom}_R(M,N)$ also extends $f$, then $\phi\left(\sum\limits_{b\in F}r_bb\right)=\sum\limits_{b\in F}r_bf(b)=\widetilde{f}\left(\sum\limits_{b\in F}r_bb\right)$, hence $\phi=\widetilde{f}$.\\
\hlt{Structural Commentary}: A free module is controlled entirely by its basis values.\\
\hlt{Key Insight}: Freeness is the universal permission to choose images of basis elements arbitrarily.
\end{proof}
\end{theorem}

\begin{theorem} \label{thm:free-module-rank-well-defined-commutative} \hlt{Rank Is Well Defined fo r Free Modules}\\
Let $A$ benonzero commutative ring with identity, $M$ be free $A$-module, $\operatorname{Bas}_A(M)$ denote set of $A$-bases of $M$.
\begin{enumerate}[label=(\roman*)]
\setlength{\itemsep}{0pt}
\item $B,B'\in\operatorname{Bas}_A(M)\Longrightarrow\abs{B}=\abs{B'}$.
\item $B\in\operatorname{Bas}_A(M)$, $S\subseteq M$, and $\langle\!\langle S\rangle\!\rangle_A=M\Longrightarrow\abs{B}\leq\abs{S}$.
\end{enumerate}
\begin{proof}
\begin{enumerate}[label=(\roman*)]
\setlength{\itemsep}{0pt}
\item Choose a maximal ideal $\mathcal{I}\lhd A$ by Theorem \ref{thm:maximal-ideals-exist}.\\
Put $K\defeq A/\mathcal{I}$, a field by Theorem \ref{thm:quotient-by-maximal-ideal-is-field}, $\mathcal{I}M\defeq\langle\!\langle im:i\in\mathcal{I},\ m\in M\rangle\!\rangle_A$, and $V\defeq M/\mathcal{I}M$.\\
Since $\mathcal{I}V=0$, $V$ is a $K$-vector space via $(a+\mathcal{I})(m+\mathcal{I}M)\defeq am+\mathcal{I}M$.
\item Let $B\in\operatorname{Bas}_A(M)$ and set $\overline{B}\defeq\{\overline{b}\defeq b+\mathcal{I}M:b\in B\}$.\\
In $B$-coordinates, $\mathcal{I}M=\{\sum\limits_{b\in F}i_bb:F\subseteq_{\operatorname{fin}}B,\ i_b\in\mathcal{I}\}$; hence $\overline{B}$ spans $V$.\\
If $\sum\limits_{j=1}^n(a_j+\mathcal{I})\overline{b_j}=0$ with distinct $b_j\in B$, then $\sum\limits_{j=1}^na_jb_j\in\mathcal{I}M$, so $a_j\in\mathcal{I}$ $\forall j$; hence $\overline{B}$ is $K$-independent.\\
If $\overline b=\overline{b'}$ with $b\neq b'$, then $b-b'\in\mathcal{I}M$, forcing the $B$-coefficient $1\in\mathcal{I}$, impossible. Thus $B\xrightarrow{\sim}\overline B$.
\item Therefore $\abs{B}=\abs{\overline B}=\dim_K V$ for every $B\in\operatorname{Bas}_A(M)$, so (i) follows from Theorem \ref{thm:arbitrary-bases-same-cardinality}.\\
If $\langle\!\langle S\rangle\!\rangle_A=M$, then $\overline S\defeq\{s+\mathcal{I}M:s\in S\}$ spans $V$, so $\abs{B}=\dim_K V\leq\abs{\overline S}\leq\abs S$.
\end{enumerate}
\hlt{Structural Commentary}: Modulo a maximal ideal, free module becomes a vector space with same basis size.\\
\hlt{Key Insight}: The field quotient $A/\mathcal{I}$ lets vector-space dimension certify module rank.
\end{proof}
\end{theorem}

\begin{definition} \label{def:rank-of-free-module} \hlt{Rank of a Free Module}\\
Let $A$ be a nonzero commutative ring with identity. The \hlt{rank} of a free
$A$-module $M$, denoted $\operatorname{rk}_A(M)$, is the cardinality of any
basis of $M$. The zero module has rank $0$.
\end{definition}

\begin{example} \label{ex:noncommutative-rank-failure-endomorphism-ring} \hlt{Why Commutativity Matters for Rank}\\
Let $V=\bigoplus_{k\geq0}\mathbb{F}e_k$ and $E=\LL(V)$. As a left $E$-module, $E=E\operatorname{id}_V$, so $\{\operatorname{id}_V\}$ is a basis.\\
For each $n\geq1$, define $\beta_s\in E$, $0\leq s<n$, by $\beta_s(e_{kn+t})=
\begin{cases}
e_k,&t=s,\\
0,&t\neq s,
\end{cases}$, where $k\geq0$, $0\leq t<n$.\\
If $\sum\limits_{s=0}^{n-1}\alpha_s\beta_s=0$, then evaluating at $e_{kn+t}$ gives $\alpha_t(e_k)=0$ for all $k,t$, hence $\alpha_t=0$ for all $t$.\\
Conversely, for $h\in E$, define $\alpha_s(e_k)\defeq h(e_{kn+s})$; then $h=\sum\limits_{s=0}^{n-1}\alpha_s\beta_s$.\\
Thus $\{\beta_0,\ldots,\beta_{n-1}\}$ is an $E$-basis of $E$, so $E$ has bases of all positive finite cardinalities.\\
\hlt{Structural Commentary}: Noncommutative scalar rings can make a free module have incompatible basis sizes.\\
\hlt{Key Insight}: The rank theorem is a commutative-ring theorem, not a formal consequence of freeness alone.
\end{example}

\begin{theorem} \label{thm:finite-support-functions-free-module} \hlt{Finite-Support Coordinate Modules}\\
Let $A$ be a commutative ring with identity and let $B$ be a set. Put $A^{(B)}\defeq\{f\in A^B:\abs{\operatorname{supp}(f)}<\infty\}$.\\
For $b\in B$, let $\delta_b\in A^{(B)}$ be given by $\delta_b(b)=1$ and $\delta_b(x)=0$ for $x\neq b$.\\
Then $A^{(B)}=\bigoplus\limits_{b\in B}A\delta_b$, so $\{\delta_b:b\in B\}$ is an $A$-basis. If $A\neq0$, then $\operatorname{rk}_A(A^{(B)})=\abs{B}$.
\begin{proof}
For $f,g\in A^{(B)}$ and $a\in A$, $\operatorname{supp}(f+g)\subseteq\operatorname{supp}(f)\cup\operatorname{supp}(g)$ and $\operatorname{supp}(af)\subseteq\operatorname{supp}(f)$, hence $A^{(B)}\leq A^B$.\\
For $f\in A^{(B)}$, $f=\sum\limits_{b\in\operatorname{supp}(f)}f(b)\delta_b$, so the $\delta_b$ span.
If $\sum\limits_{j=1}^{m}a_j\delta_{b_j}=0$ with $b_1,\ldots,b_m$ distinct, then evaluation at $b_i$ gives $a_i=0$ for every $i$. Thus the $\delta_b$ are independent, and the rank statement follows from the displayed basis.\\
\hlt{Structural Commentary}: Free modules are finite-support coordinate modules on their bases.\\
\hlt{Key Insight}: Finite support is the module-theoretic version of finite linear combination.
\end{proof}
\end{theorem}

\begin{proposition} \label{prop:free-module-coordinate-model-and-isomorphism} \hlt{Free Modules Are Coordinate Modules}\\
Let $A$ be a commutative ring with identity and let $M$ be free on a basis $B$.\\
Then the coordinate map $\Phi_B:M\to A^{(B)}$ is an isomorphism.\\
If $A\neq0$ and $M,N$ are free $A$-modules, then $M\cong N\Longleftrightarrow \operatorname{rk}_A(M)=\operatorname{rk}_A(N)$.
\begin{proof} For $m=\sum\limits_{b\in F}a_bb$, set $\Phi_B(m)=f_m$, with $f_m(b)=a_b$ for $b\in F$ and $f_m(b)=0$ for $b\notin F$.\\
Unique coordinates give $\Phi_B:M\xrightarrow{\sim}A^{(B)}$, with inverse $f\mapsto\sum\limits_{b\in\operatorname{supp}(f)}f(b)b$.
\begin{enumerate}[label=\roman*.]
\setlength{\itemsep}{0pt}
\item[$(\Rightarrow)$] If $\psi:M\xrightarrow{\sim}N$ and $C$ is a basis of $M$, then $\psi(C)$ is a basis of $N$, so $\operatorname{rk}_A(M)=\abs{C}=\abs{\psi(C)}=\operatorname{rk}_A(N)$.
\item[$(\Leftarrow)$] If $\operatorname{rk}_A(M)=\operatorname{rk}_A(N)$, choose bases $C,D$ with $\abs{C}=\abs{D}$.\\
Any bijection $C\to D$ extends, by Theorem \ref{thm:free-module-map-extension}, to $M\xrightarrow{\sim}N$.
\end{enumerate}
\noindent
\hlt{Structural Commentary}: A basis identifies a free module with a standard coordinate model.\\
\hlt{Key Insight}: Over a nonzero commutative ring, rank classifies free modules up to isomorphism.
\end{proof}
\end{proposition}

\begin{theorem} \label{thm:independent-sets-bounded-by-free-rank-domain} \hlt{Independent Sets Are Bounded by Rank over a Domain}\\
Let $A$ be integral domain with field of quotients $Q$, $M$ free over $A$, $S\subseteq M$ $A$-independent. Then $\abs{S}\leq\operatorname{rk}_A(M)$.
\begin{proof}
Choose an $A$-basis $B$ and identify $M=A^{(B)}\subseteq Q^{(B)}$.\\
If $\sum\limits_{i=1}^{n}q_is_i=0$ with $q_i\in Q$ and distinct $s_i\in S$, choose $d\in A\setminus\{0\}$ with $dq_i\in A$ for all $i$.\\
Then $\sum\limits_{i=1}^{n}(dq_i)s_i=0$ in $M$, so $dq_i=0$ by $A$-independence, hence $q_i=0$ for all $i$.\\
Thus $S$ is $Q$-independent in $Q^{(B)}$, and $\abs{S}\leq\dim_Q Q^{(B)}=\abs{B}=\operatorname{rk}_A(M)$.\\
\hlt{Structural Commentary}: The field of quotients lets ordinary dimension bound module independence.\\
\hlt{Key Insight}: Domains allow denominator clearing without creating false zero relations.
\end{proof}
\end{theorem}

\begin{theorem} \label{thm:epimorphism-onto-free-module-splits} \hlt{Free Quotients Split}\\
Let $A$ be a commutative ring with identity.
\begin{enumerate}[label=(\roman*)]
\setlength{\itemsep}{0pt}
\item If $\sigma:M\twoheadrightarrow F$ and $F$ is free, then $\exists N\leq M$ such that $M=\ker(\sigma)\oplus N$ and $N\cong F$.
\item If $S\leq M$ and $M/S$ is free, then $\exists T\leq M$ such that $M=S\oplus T$ and $T\cong M/S$; hence $M\cong S\oplus M/S$.
\item If $A\neq0$ and $M,S,M/S$ are free, then $\operatorname{rk}_A(M)=\operatorname{rk}_A(S)+\operatorname{rk}_A(M/S)$, with cardinal addition; in finite rank, $\operatorname{rk}_A(M/S)=\operatorname{rk}_A(M)-\operatorname{rk}_A(S)$.
\end{enumerate}
\begin{proof}
\begin{enumerate}[label=(\roman*)]
\setlength{\itemsep}{0pt}
\item Let $B$ be a basis of $F$. For $b\in B$, choose $m_b\in M$ with $\sigma(m_b)=b$. By Theorem \ref{thm:free-module-map-extension}, $b\mapsto m_b$ extends to $\tau:F\to M$. Since $\sigma\tau=\operatorname{id}_F$, Theorem
\ref{thm:module-one-sided-inverses-and-complements} gives $M=\ker(\sigma)\oplus\operatorname{im}(\tau)$ and $\operatorname{im}(\tau)\cong F$.
\item Apply (i) to $\pi:M\twoheadrightarrow M/S$. Then $M=S\oplus T$ for some $T\cong M/S$, so $M\cong S\oplus M/S$.
\item From (ii), $M=S\oplus T$ with $T\cong M/S$. If $B_S,B_T$ are bases of $S,T$, then $B_S\sqcup B_T$ is a basis of $M$, hence $\operatorname{rk}_A(M)=\abs{B_S}+\abs{B_T} =\operatorname{rk}_A(S)+\operatorname{rk}_A(M/S)$. The finite formula follows.
\end{enumerate}
\noindent
\hlt{Structural Commentary}: Free modules are projective: maps onto them have linear sections.\\
\hlt{Key Insight}: Choosing representatives for basis vectors of the quotient is enough to split the quotient map.
\end{proof}
\end{theorem}

\subsection{Noetherian Modules}

\begin{definition} \label{def:noetherian-module-ring-acc} \hlt{Noetherian Modules and Rings}\\
Let $A$ be a commutative ring with identity.
\begin{enumerate}[label=(\roman*)]
\setlength{\itemsep}{0pt}
\item An $A$-module $M$ satisfies \hlt{ACC} if every chain $M_1\leq M_2\leq\cdots$ of submodules satisfies\\
$\exists n_0, \forall n\geq n_0, M_n=M_{n_0}$. Such $M$ is \hlt{Noetherian}.
\item The ring $A$ is \hlt{Noetherian} if the regular $A$-module $A$ is Noetherian,\\ equivalently if every ideal-chain $I_1\subseteq I_2\subseteq\cdots$ stabilizes.
\end{enumerate}
\end{definition}

\begin{theorem} \label{thm:noetherian-iff-submodules-finitely-generated} \hlt{Noetherian Means Finitely Generated Submodules}\\
Let $M$ be an $A$-module.
\begin{enumerate}[label=(\roman*)]
\setlength{\itemsep}{0pt}
\item $M$ is Noetherian $\Longleftrightarrow$ every $N\leq M$ is f.g.
\item If $A$ is commutative with identity, then $A$ is Noetherian
$\Longleftrightarrow$ every $\mathcal{I}\lhd A$ is f.g.
\end{enumerate}
\begin{proof}
\begin{enumerate}[label=(\roman*).]
\setlength{\itemsep}{0pt}
\item[$(i)$]($\Rightarrow$) Let $N\leq M$. If $N$ is not f.g., choose $u_{n+1}\in N\setminus\langle\!\langle u_1,\ldots,u_n\rangle\!\rangle_A$ inductively.\\
Then $\langle\!\langle u_1\rangle\!\rangle_A< \langle\!\langle u_1,u_2\rangle\!\rangle_A<\cdots$, contradicting ACC.
\item[$(i)$]($\Leftarrow$) Let $M_1\leq M_2\leq\cdots$ and put $U=\bigcup_{n\geq1}M_n$. Then $U\leq M$, so $U=\langle\!\langle u_1,\ldots,u_r\rangle\!\rangle_A$.\\
Choose $k$ with $u_1,\ldots,u_r\in M_k$. Hence $U=M_k$ and $M_n=M_k$ for all $n\geq k$.
\item[$(ii)$] Apply (i) to the regular $A$-module $A$: its submodules are exactly the ideals of $A$.
\end{enumerate}
\noindent
\hlt{Structural Commentary}: ACC and finite generation are two faces of the same finiteness condition.\\
\hlt{Key Insight}: Failure of finite generation builds an infinite increasing chain one new generator at a time.
\end{proof}
\end{theorem}

\begin{theorem} \label{thm:finitely-generated-modules-over-noetherian-rings} \hlt{Finitely Generated Modules over Noetherian Rings}\\
Let $A$ be a commutative ring with identity. Then $A$ is Noetherian $\Longleftrightarrow$ every f.g. $A$-module is Noetherian.
\begin{proof}
\begin{enumerate}[label=\roman*.]
\setlength{\itemsep}{0pt}
\item[$(\Rightarrow)$] Assume $A$ Noetherian. First prove by induction on $n$ that every $S\leq A^n$ is f.g.\\
The case $n=1$ is Theorem \ref{thm:noetherian-iff-submodules-finitely-generated}. For $n>1$, let $\pi:S\to A$ be last-coordinate projection.\\
Then $\pi(S)\lhd A$ is f.g. and $\ker(\pi)\leq A^{n-1}$ is f.g. by induction. Lifts of generators of $\pi(S)$, together with generators of $\ker(\pi)$, generate $S$. Thus $A^n$ is Noetherian. If $M$ is f.g., choose $A^n\twoheadrightarrow M$; for $N\leq M$, its preimage in $A^n$ is f.g., hence $N$ is f.g. Therefore $M$ is Noetherian.
\item[$(\Leftarrow)$] Since $A=\langle\!\langle1\rangle\!\rangle_A$, the regular $A$-module $A$ is f.g.; by hypothesis it is Noetherian. Hence $A$ is Noetherian.
\end{enumerate}
\noindent
\hlt{Structural Commentary}: Noetherianity passes from rings to finite free modules to quotients.\\
\hlt{Key Insight}: Finitely generated modules are quotients of finite free modules.
\end{proof}
\end{theorem}

\subsection{Smith Normal Form}

\begin{definition} \label{def:standing-pid-notation} \hlt{Standing PID Notation}\\
In the PID-structure results, $R$ denotes a principal ideal domain unless otherwise stated.\\
For $r\in R$, write $\langle r\rangle=Rr$; for an $R$-module $M$ and $v\in M$, write $\langle\!\langle v\rangle\!\rangle_R=Rv$.
\end{definition}

\begin{perspective} \label{per:pid-modules-scalar-arithmetic-controls-structure} \hlt{PID Arithmetic Controls Module Structure}\\
Fields have $F\setminus\{0\}\subseteq F^\times$. PIDs replace invertibility by factorization, gcds, and lcms.\\
This divisibility data classifies finitely generated modules.
\end{perspective}

\begin{definition} \label{def:module-order-over-pid} \hlt{Order of a Submodule or Element}\\
Let $M$ be an $R$-module and $N\leq M$. Since $\operatorname{Ann}_R(N)=\{r\in R:rN=0\}=\langle\alpha\rangle$, any such $\alpha$ is an \hlt{order} of $N$; write $o(N)=\alpha$, up to associates. For $v\in M$, set $o(v)\defeq o(\langle\!\langle v\rangle\!\rangle_R)$.
\end{definition}

\begin{proposition} \label{prop:orders-direct-sum-lcm} \hlt{Orders and Direct Sums}\\
Let $M=A\oplus B$. If $o(A)=\alpha$ and $o(B)=\beta$, then $o(M)\sim\operatorname{lcm}(\alpha,\beta)$.\\
More generally, $o(\bigoplus\limits_{i=1}^{n}M_i)\sim\operatorname{lcm}(o(M_1),\ldots,o(M_n))$.
\begin{proof}
$r(A\oplus B)=0\Longleftrightarrow rA=0$ and $rB=0$.\\
Hence $\operatorname{Ann}_R(A\oplus B)=\operatorname{Ann}_R(A)\cap\operatorname{Ann}_R(B) =\langle\alpha\rangle\cap\langle\beta\rangle =\langle\operatorname{lcm}(\alpha,\beta)\rangle$.\\
Thus $o(A\oplus B)\sim\operatorname{lcm}(\alpha,\beta)$, and induction gives the finite case.\\
\hlt{Structural Commentary}: Direct sums intersect annihilator conditions.\\
\hlt{Key Insight}: Lcms appear because one scalar must kill every summand at once.
\end{proof}
\end{proposition}

\begin{theorem} \label{thm:cyclic-modules-over-pid} \hlt{Cyclic Modules over a PID}\\
Let $M$ be an $R$-module and let $v\in M$.
\begin{enumerate}[label=(\roman*)]
\setlength{\itemsep}{0pt}
\item If $\operatorname{Ann}_R(v)=\langle\alpha\rangle$, then $R/\langle\alpha\rangle\xrightarrow{\sim}\langle\!\langle v\rangle\!\rangle_R$, $r+\langle\alpha\rangle\mapsto rv$.
\item If $C=\langle\!\langle v\rangle\!\rangle_R$ and $S\leq C$, then $S=\langle\!\langle sv\rangle\!\rangle_R$ for some $s\in R$.
\item If $o(v)=\alpha$ and $\beta\in R$, then $o(\beta v)\sim\alpha/\gcd(\alpha,\beta)$, and $\langle\!\langle\beta v\rangle\!\rangle_R=\langle\!\langle v\rangle\!\rangle_R \Longleftrightarrow \gcd(\alpha,\beta)\sim1$.
\end{enumerate}
\begin{proof}
\begin{enumerate}[label=(\roman*)]
\setlength{\itemsep}{0pt}
\item The map $\phi:R\to\langle\!\langle v\rangle\!\rangle_R$, $\phi(r)=rv$, is onto and $\ker\phi=\operatorname{Ann}_R(v)=\langle\alpha\rangle$. Hence $R/\langle\alpha\rangle\xrightarrow{\sim}\langle\!\langle v\rangle\!\rangle_R$.
\item For $S\leq C$, set $I=\{r\in R:rv\in S\}\lhd R$.\\
Since $R$ is a PID, $I=\langle s\rangle$ for some $s\in R$, and $S=\{rv:r\in I\}=\langle\!\langle sv\rangle\!\rangle_R$.
\item Put $d=\gcd(\alpha,\beta)$, $\alpha=d\alpha'$, and $\beta=d\beta'$ with $\gcd(\alpha',\beta')\sim1$.\\
For $r\in R$, $r\beta v=0\Longleftrightarrow\alpha\mid r\beta \Longleftrightarrow\alpha'\mid r\beta'\Longleftrightarrow\alpha'\mid r$, the last step by Bezout.\\
Thus $\operatorname{Ann}_R(\beta v)=\langle\alpha'\rangle =\langle\alpha/\gcd(\alpha,\beta)\rangle$. Finally,
$\langle\!\langle\beta v\rangle\!\rangle_R=\langle\!\langle v\rangle\!\rangle_R$ $\Longleftrightarrow(\exists c\in R)(v=c\beta v)$\\
$\Longleftrightarrow(\exists c\in R)(\alpha\mid1-c\beta) \Longleftrightarrow\gcd(\alpha,\beta)\sim1$.
\end{enumerate}
\noindent
\hlt{Structural Commentary}: A cyclic module is the base ring modulo the first relation satisfied by its generator.\\
\hlt{Key Insight}: Multiplying a cyclic generator divides its order by the gcd of the multiplier and the old order.
\end{proof}
\end{theorem}

\begin{theorem} \label{thm:cyclic-composition-coprime-orders} \hlt{Composing and Decomposing Cyclic Modules}\\
Let $M$ be an $R$-module.
\begin{enumerate}[label=(\roman*)]
\setlength{\itemsep}{0pt}
\item If $o(u_i)=\alpha_i\neq0$ and $\gcd(\alpha_i,\alpha_j)\sim1$ for $i\neq j$, then $o(\sum\limits_{i=1}^{n}u_i)\sim\prod_{i=1}^{n}\alpha_i$ and $\bigoplus\limits_{i=1}^{n}\langle\!\langle u_i\rangle\!\rangle_R =\langle\!\langle\sum\limits_{i=1}^{n}u_i\rangle\!\rangle_R$.
\item If $o(v)=\alpha=\prod\limits_{i=1}^{n}\alpha_i\neq0$ and $\gcd(\alpha_i,\alpha_j)\sim1$ for $i\neq j$,\\
then $\exists v_1,\ldots,v_n\in M$ such that $v=\sum\limits_{i=1}^{n}v_i$, $o(v_i)\sim\alpha_i$, and $\langle\!\langle v\rangle\!\rangle_R =\bigoplus\limits_{i=1}^{n}\langle\!\langle v_i\rangle\!\rangle_R$.
\end{enumerate}
\begin{proof}
\begin{enumerate}[label=(\roman*)]
\setlength{\itemsep}{0pt}
\item It suffices to prove $n=2$ and iterate. Let $o(u)=\alpha$, $o(w)=\beta$, and $\gcd(\alpha,\beta)\sim1$.\\
If $x\in\langle\!\langle u\rangle\!\rangle_R\cap\langle\!\langle w\rangle\!\rangle_R$, then $\alpha x=\beta x=0$, so Bezout gives $x=0$.  If $r(u+w)=0$, then $ru=-rw$ lies in this intersection; hence $ru=rw=0$, so $\alpha\mid r$ and $\beta\mid r$, hence $\alpha\beta\mid r$. Since $\alpha\beta(u+w)=0$, $o(u+w)\sim\alpha\beta$. Choose $a,b\in R$ with $a\alpha+b\beta=1$. Then $u=b\beta(u+w)$ and $w=a\alpha(u+w)$, so $\langle\!\langle u\rangle\!\rangle_R\oplus \langle\!\langle w\rangle\!\rangle_R=\langle\!\langle u+w\rangle\!\rangle_R$.
\item Put $\beta_i=\alpha/\alpha_i$. Since $\gcd(\beta_1,\ldots,\beta_n)\sim1$, choose $c_i\in R$ with $\sum\limits_{i=1}^{n}c_i\beta_i=1$, and set $v_i=c_i\beta_i v$.\\
Then $v=\sum\limits_{i=1}^{n}v_i$. Also $\gcd(c_i,\alpha_i)\sim1$ and $\gcd(\alpha,\beta_i)\sim\beta_i$, so Theorem \ref{thm:cyclic-modules-over-pid} gives $o(v_i)\sim\alpha_i$.\\
Part (i) gives $\langle\!\langle v\rangle\!\rangle_R =\bigoplus\limits_{i=1}^{n}\langle\!\langle v_i\rangle\!\rangle_R$.
\end{enumerate}
\noindent
\hlt{Structural Commentary}: Relatively prime cyclic pieces merge into one cyclic piece.\\
\hlt{Key Insight}: Bezout identities recover each component from the single sum.
\end{proof}
\end{theorem}

\begin{perspective} \label{per:pid-noetherian-specialization} \hlt{Noetherian Modules over a PID}\\
Over a PID, $N\leq R^n$ implies $N$ free with $\operatorname{rk}_R(N)\leq n$; Noetherianity is the f.g. shadow.
\end{perspective}

\begin{theorem} \label{thm:submodule-of-finite-free-pid-is-free} \hlt{Submodules of Finite Free PID Modules}\\
Let $F$ be free over $R$, $\operatorname{rk}_R(F)=n<\infty$, and $N\leq F$. Then $\exists m\leq n$ such that $N\cong R^m$.\\
Consequently, for $\mu_R(M)\defeq\min\{\abs{X}:\langle\!\langle X\rangle\!\rangle_R=M\}$, $\mu_R(M)\leq n$ implies $S\leq M\Longrightarrow \mu_R(S)\leq n$.
\begin{proof}
Identify $F\simeq R^n$. For $1\leq k\leq n$, set $I_k\defeq\{x_k:x=(x_1,\ldots,x_k,0,\ldots,0)\in N\}\lhd R$ and write $I_k=\langle r_k\rangle$..\\
If $r_k\neq0$, choose $u_k\in N\cap(R^k\times\{0\}^{n-k})$ with $(u_k)_k=r_k$. Let $K=\{k:r_k\neq0\}$.\\
If $\sum\limits_{k\in K}c_ku_k=0$ and $j=\max\{k:c_k\neq0\}$, then the $j$th coordinate gives $c_jr_j=0$, hence $c_j=0$, contradiction.\\
Thus $\{u_k:k\in K\}$ is $R$-independent. Set $\ell(0)=0$ and $\ell(x)=\max\operatorname{supp}(x)$ for $0\neq x\in N$.\\
If $j=\ell(x)>0$, then $x_j\in I_j=\langle r_j\rangle$, say $x_j=cr_j$, and $\ell(x-cu_j)<j$.\\
Induction on $\ell(x)$ gives $N=\bigoplus\limits_{k\in K}Ru_k\cong R^{\abs{K}}$, so $\operatorname{rk}_R(N)=\abs{K}\leq n$.\\
If $\mu_R(M)\leq n$, choose $\pi:R^n\twoheadrightarrow M$. For $S\leq M$, $P\defeq\pi^{-1}(S)\leq R^n$ has $P=\langle\!\langle v_1,\ldots,v_m\rangle\!\rangle_R$ with $m\leq n$.\\
Hence $S=\langle\!\langle \pi(v_1),\ldots,\pi(v_m)\rangle\!\rangle_R$ and $\mu_R(S)\leq m\leq n$.\\
\hlt{Structural Commentary}: PIDs let one choose pivot generators in a submodule of $R^n$.\\
\hlt{Key Insight}: Pivot ideals $I_k$ yield triangular free generators.
\end{proof}
\end{theorem}

\begin{corollary} \label{cor:finitely-generated-pid-modules-noetherian} \hlt{Finitely Generated PID Modules Are Noetherian}\\
If $\mu_R(M)<\infty$, then $M$ is Noetherian.
\begin{proof} Let $n=\mu_R(M)$. By Theorem \ref{thm:submodule-of-finite-free-pid-is-free}, $N\leq M\Longrightarrow \mu_R(N)\leq n<\infty$.\\
Thus every submodule of $M$ is f.g.; Theorem \ref{thm:noetherian-iff-submodules-finitely-generated} gives $M$ Noetherian.\\
\hlt{Structural Commentary}: Finite generation of every submodule prevents infinite strict growth.\\
\hlt{Key Insight}: Generators of the union appear at a finite stage.
\end{proof}
\end{corollary}

\begin{definition} \label{def:rank-finite-free-module-pid} \hlt{Rank of a Finite Free Module}\\
If $F$ is a finite free $R$-module and $F\cong R^n$, the integer $n$ is the \hlt{rank} of $F$, denoted $\operatorname{rk}(F)$.
\end{definition}

\begin{proposition} \label{prop:finite-free-rank-well-defined-domain} \hlt{Finite Free Rank Is Well-Defined over a Domain}\\
Let $R$ be an integral domain and $m,n\geq0$. Then $R^m\cong_R R^n\Longrightarrow m=n$.
\begin{proof}
Let $K=\operatorname{Frac}(R)$ and $\phi:R^m\xrightarrow{\sim}R^n$. Then $K\otimes_R\phi:K\otimes_R R^m\xrightarrow{\sim}K\otimes_R R^n$, i.e. $K^m\xrightarrow{\sim}K^n$.\\
Hence $m=\dim_KK^m=\dim_KK^n=n$.\\
\hlt{Structural Commentary}: Passing to the fraction field recovers ordinary dimension.\\
\hlt{Key Insight}: Rank over a domain is detected after denominators are allowed.
\end{proof}
\end{proposition}

\begin{proposition} \label{prop:spanning-set-size-rank-basis-pid} \hlt{A Full-Size Spanning Set Is a Basis}\\
Let $\operatorname{rk}_R(F)=n<\infty$. If $F=\langle\!\langle s_1,\ldots,s_n\rangle\!\rangle_R$, then $(s_1,\ldots,s_n)$ is an $R$-basis of $F$.
\begin{proof}
Choose a basis $(b_i)_{i=1}^n$ and define $\tau\in\operatorname{End}_R(F)$ by $\tau(b_i)=s_i$. Then $\tau$ is onto.\\
Since $F$ is free, $\exists\rho:F\to F$ with $\tau\rho=\operatorname{id}_F$; hence $F=\ker\tau\oplus\operatorname{im}\rho$ and $\operatorname{im}\rho\cong F$.\\
By Theorem \ref{thm:submodule-of-finite-free-pid-is-free}, $\ker\tau\cong R^r$ for some $r\leq n$. Thus $R^n\cong F\cong\ker\tau\oplus F\cong R^{r+n}$, so $r=0$ by Proposition
\ref{prop:finite-free-rank-well-defined-domain}. Hence $\ker\tau=0$, $\tau$ is an isomorphism, and $(s_i)=\tau(b_i)$ is a basis.\\
\hlt{Structural Commentary}: Surjective finite-free endomorphisms over a PID are injective.\\
\hlt{Key Insight}: Rank leaves no room for a split kernel.
\end{proof}
\end{proposition}

\begin{perspective} \label{per:smith-normal-form-and-module-presentations} \hlt{Smith Form as Diagonalized Relations}\\
For a presentation $R^n\xrightarrow{A}R^m\to M\to0$, Smith form replaces $A$ by a diagonal relation matrix, decoupling finite relations into cyclic constraints.
\end{perspective}

\begin{theorem} \label{thm:smith-normal-form-over-pid} \hlt{Smith Normal Form over a PID}\\
Let $A\in\operatorname{Mat}(m,n,R)$, viewed as $R^n\to R^m$.\\
$\exists$ invertible $P\in\operatorname{Mat}(m,m,R)$ and $Q\in\operatorname{Mat}(n,n,R)$ such that $PAQ=\operatorname{diag}(d_1,\ldots,d_r,0,\ldots,0)$,\\
where $d_i\in R\setminus\{0\}$, $d_1\mid\cdots\mid d_r$, and $r=\operatorname{rk}_R(\operatorname{im}A)$. The ideals $\langle d_i\rangle$ are unique.
\begin{proof}
\begin{enumerate}[label=(\roman*)]
\setlength{\itemsep}{0pt}
\item \hlt{Existence}: If $A=0$, done. Set $I(A)=\langle a_{ij}\rangle=\langle d\rangle$. Elementary operations move a Bezout combination equal to $d$ to position $(1,1)$. If a pivot $p\nmid a$, Bezout on $(p,a)$ replaces $p$ by $\gcd(p,a)$, so $\langle p\rangle$ strictly increases; ACC in a PID forces termination. Then $p$ divides every entry, so clear its row and column.\\
Induct on the complementary submatrix. Finally replace each bad adjacent pair $(a,b)$ by $(\gcd(a,b),\operatorname{lcm}(a,b))$ until $d_1\mid\cdots\mid d_r$.
\item \hlt{Uniqueness}: Let $\Delta_k(A)$ be the ideal of $k\times k$ minors and $\Delta_0(A)=R$. Matrix equivalence preserves each $\Delta_k$. In Smith form, $\Delta_k(A)=\langle d_1\cdots d_k\rangle$ for $1\leq k\leq r$, and $\Delta_k(A)=0$ for $k>r$. Since $R$ is a domain, $(\Delta_k(A):\Delta_{k-1}(A))=\langle d_k\rangle$, so $d_k$ is determined up to associates.
\end{enumerate}
\noindent
\hlt{Structural Commentary}: Smith form preserves rank plus divisibility data.\\
\hlt{Key Insight}: Determinantal ideals encode invariant factors.
\end{proof}
\end{theorem}

\begin{theorem} \label{thm:adapted-basis-submodule-free-pid} \hlt{Adapted Bases for Submodules of Free Modules}\\
Let $N\leq F$, $\operatorname{rk}_R(F)=n$, and $\operatorname{rk}_R(N)=k$.\\
Then $\exists(v_i)_{i=1}^n\in\operatorname{Bas}_R(F)$ and $r_1,\ldots,r_k\in R\setminus\{0\}$ such that $r_1\mid\cdots\mid r_k$ and $(r_iv_i)_{i=1}^k\in\operatorname{Bas}_R(N)$.
\begin{proof}
Choose bases of $N$ and $F$, and let $A$ represent $\iota:N\hookrightarrow F$.\\
Smith form gives $PAQ=\operatorname{diag}(r_1,\ldots,r_k,0,\ldots,0)$ with $r_i\neq0$ and $r_1\mid\cdots\mid r_k$, since $\iota$ is injective.\\
The row operations give a basis $(v_i)_{i=1}^n$ of $F$ and the column operations give a basis $(w_i)_{i=1}^k$ of $N$ with $\iota(w_i)=r_iv_i$. Hence $(r_iv_i)_{i=1}^k=(w_i)_{i=1}^k$ is a basis of $N$.\\
\hlt{Structural Commentary}: Inclusions diagonalize after choosing bases.\\
\hlt{Key Insight}: Diagonal entries are the required rescalings.
\end{proof}
\end{theorem}

\begin{corollary} \label{cor:smith-normal-form-presents-pid-modules} \hlt{Smith Form Presents Finitely Generated PID Modules}\\
Let $R^n\xrightarrow{A}R^m\to M\to0$ be a finite presentation.\\
If $A\sim\operatorname{diag}(d_1,\ldots,d_r,0,\ldots,0)$ in Smith form. $I=\{i:1\leq i\leq r,\ d_i\notin R^\times\}$, then $M\cong R^{m-r}\oplus\bigoplus_{i\in I}R/\langle d_i\rangle$.
\begin{proof}
Equivalence $A\sim PAQ$ changes bases in $R^m,R^n$, so $\operatorname{coker}A\cong\operatorname{coker}(PAQ)$.\\
For the diagonal map, $d_i\in R^\times$ gives $0$, $d_i\notin R^\times$ gives $R/\langle d_i\rangle$, and each zero row gives $R$.\\
\hlt{Structural Commentary}: Smith form turns presentations into direct sums.\\
\hlt{Key Insight}: Diagonal relations create cyclic torsion summands, and zero rows create free summands.
\end{proof}
\end{corollary}

\begin{method} \label{meth:smith-form-module-presentation} \hlt{Reading a Module from Smith Form}\\
For finite presentation $R^n\xrightarrow{A}R^m\to M\to0$ over PID, put $A\sim D=\operatorname{diag}(d_1,\ldots,d_r,0,\ldots,0)$, $d_1\mid\cdots\mid d_r$:
\begin{enumerate}[label=(\roman*)]
\setlength{\itemsep}{0pt}
\item $d_i\in R^\times\Longrightarrow$ no summand;
\item $d_i\notin R^\times\Longrightarrow R/\langle d_i\rangle$;
\item each zero target coordinate contributes one $R$;
\item factor nonunit $d_i$ into prime powers for elementary divisors.
\end{enumerate}
\hlt{Key Insight}: Smith form changes arbitrary relations into independent one-coordinate relations.
\end{method}

\begin{example} \label{ex:smith-normal-form-abelian-group-decomposition} \hlt{Smith Form for a Finite Abelian Group}\\
Let $M=\mathbb{Z}^2/\langle\!\langle(6,0),(0,10),(2,4)\rangle\!\rangle_{\mathbb{Z}}$, with relation matrix $A=\begin{bmatrix}6&0&2\\0&10&4\end{bmatrix}$. Then $\Delta_1(A)=\langle6,10,2,4\rangle=\langle2\rangle$ and $\Delta_2(A)=\langle60,24,-20\rangle=\langle4\rangle$. Hence $d_1\sim2$, $d_1d_2\sim4$, so $d_2\sim2$ and $M\cong\mathbb{Z}/2\mathbb{Z}\oplus\mathbb{Z}/2\mathbb{Z}$.\\
\hlt{Key Insight}: Determinantal ideals read off cyclic summands.
\end{example}

\begin{example} \label{ex:elementary-divisors-recombine-to-invariant-factors} \hlt{Recombining Elementary Divisors}\\
Let $G=\Z/12\Z\oplus\Z/18\Z$. Then $\operatorname{ElemDiv}(G)=\{2,2^2,3,3^2\}$. Recombining equal-height primary blocks gives $2\cdot3=6$ and $2^2\cdot3^2=36$; hence $G\cong\Z/6\Z\oplus\Z/36\Z$, with $6\mid36$.\\
\hlt{Structural Commentary}: Elementary divisors split by primes; invariant factors recombine them.\\
\hlt{Key Insight}: Prime-wise data equals one divisibility chain.
\end{example}

\subsection{Finitely Generated PID Modules}

\begin{theorem} \label{thm:finitely-generated-torsion-free-pid-free} \hlt{Finitely Generated Torsion-Free Modules Are Free}\\
For a f.g. $R$-module $M$, $M$ is free $\iff M$ is torsion-free.
\begin{proof}
\begin{enumerate}[label=\roman*.]
\setlength{\itemsep}{0pt}
\item[$(\Rightarrow)$] If $M\cong R^n$, then $am=0$ with $a\neq0$ forces every coordinate of $m$ to be $0$, since $R$ is a domain.
\item[$(\Leftarrow)$] Let $M$ be torsion-free and $G=\{u_1,\ldots,u_n\}$ generate $M$.\\
Choose maximal $R$-independent $S=\{v_1,\ldots,v_k\}\subseteq G$. For $w\in G\setminus S$, dependence of $S\cup\{w\}$ gives $a_w\neq0$ and $b_{w,i}\in R$ with $a_ww+\sum\limits_i b_{w,i}v_i=0$. Put $a=\prod_{w\in G\setminus S}a_w$. Then $aM\subseteq\langle\!\langle S\rangle\!\rangle_R$, so $aM$ is free by Theorem \ref{thm:submodule-of-finite-free-pid-is-free}. Since $M$ is torsion-free, $M\xrightarrow{\sim}aM$, $m\mapsto am$; hence $M$ is free.
\end{enumerate}
\noindent
\hlt{Structural Commentary}: Clear denominators into a free submodule.\\
\hlt{Key Insight}: One scalar embeds $M$ into a free sp	an.
\end{proof}
\end{theorem}

\begin{theorem} \label{thm:finitely-generated-pid-free-plus-torsion} \hlt{Free Plus Torsion Decomposition}\\
For f.g. $M$, $M=M_{\operatorname{tor}}\oplus F$ with $F$ f.g. free. Moreover $M_{\operatorname{tor}}$ is intrinsic and $F\cong M/M_{\operatorname{tor}}$.
\begin{proof} By Proposition \ref{prop:torsion-submodule-over-domain}, $T\defeq M_{\operatorname{tor}}\leq M$ and $Q\defeq M/T$ is f.g. torsion-free.\\
Thus $Q$ is free by Theorem \ref{thm:finitely-generated-torsion-free-pid-free}. Since $Q$ is free, the quotient map $\pi:M\twoheadrightarrow Q$ splits; hence $M=T\oplus\operatorname{im}\sigma$ for any section $\sigma:Q\to M$, and $\operatorname{im}\sigma\cong Q$. If $M=T'\oplus G$ with $T'$ torsion and $G$ free, then $m=t'+g\in M_{\operatorname{tor}}$ forces $g=0$, so $T'=M_{\operatorname{tor}}$. Every complement is isomorphic to $M/M_{\operatorname{tor}}$.\\
\hlt{Structural Commentary}: F.g. PID modules split into torsion plus free.\\
\hlt{Key Insight}: The quotient by torsion is free, and freeness makes the quotient map split.
\end{proof}
\end{theorem}

\begin{definition} \label{def:primary-module-over-pid} \hlt{Primary Module}\\
Let $p\in R$ be prime. A torsion $R$-module $M$ is \hlt{$p$-primary} if $o(M)\sim p^e$ for some $e\geq1$.
\end{definition}

\begin{theorem} \label{thm:primary-decomposition-pid-modules} \hlt{Primary Decomposition}\\
Let $M$ be f.g. torsion with $o(M)=\mu\sim\prod\limits_{i=1}^n p_i^{e_i}$, where the $p_i$ are distinct nonassociate primes.\\
Put $q_i\defeq p_i^{e_i}$, $\mu_i\defeq\mu/q_i$, and $M_i\defeq\mu_iM$. Then $M=\bigoplus\limits_{i=1}^nM_i$, $M_i=\{v\in M:q_iv=0\}$, and $o(M_i)\sim q_i$.\\
The decomposition is unique up to ordering.
\begin{proof}
Since $\gcd(\mu_1,\ldots,\mu_n)\sim1$, choose $a_i\in R$ with $\sum\limits_i a_i\mu_i=1$. Hence $x=\sum\limits_i a_i\mu_ix$, so $M=\sum_iM_i$.\\
Also $q_iM_i=\mu M=0$. Conversely, if $q_iv=0$, choose $a,b\in R$ with $a\mu_i+bq_i=1$; then $v=a\mu_iv\in M_i$.\\
Thus $M_i=\{v:q_iv=0\}$. If $\sum_i x_i=0$ with $x_i\in M_i$, applying $\mu_i$ gives $\mu_ix_i=0$ as $q_j\mid\mu_i$ for $j\neq i$.\\
Since $q_ix_i=0$ and $\gcd(q_i,\mu_i)\sim1$, Bezout gives $x_i=0$; hence the sum is direct.\\
Now $o(M_i)\mid q_i$ and $\operatorname{lcm}_i o(M_i)\sim o(M)\sim\prod_iq_i$, so $o(M_i)\sim q_i$.\\
Finally each summand is the intrinsic set $\{v:q_iv=0\}$, where $q_i$ is the $p_i$-part of $o(M)$.\\
\hlt{Structural Commentary}: Primary decomposition separates a torsion module prime by prime.\\
\hlt{Key Insight}: Bezout idempotents built from the coprime factors of the order project onto the primary pieces.
\end{proof}
\end{theorem}

\begin{lemma} \label{lem:p-torsion-layer-vector-space} \hlt{The First $p$-Torsion Layer}\\
Let $p\in R$ be prime and $M[p]\defeq\ker(p:M\to M)=\{v\in M:pv=0\}$.\\
Then $M[p]$ is an $R/\langle p\rangle$-vector space. If $M=S\oplus T$, then $M[p]=S[p]\oplus T[p]$.
\begin{proof} $\langle p\rangle$ is maximal by Theorem \ref{thm:pid-prime-irreducible-bezout} and Theorem \ref{thm:quotient-by-maximal-ideal-is-field}; hence $\kappa\defeq R/\langle p\rangle$ is a field. Define $(r+\langle p\rangle)v\defeq rv$.\\
If $r-r'\in\langle p\rangle$ and $v\in M[p]$, then $(r-r')v=0$, so the action is well-defined; the axioms are inherited from $M$.\\
If $v=s+t\in M[p]$ under $M=S\oplus T$, then $0=pv=ps+pt$ gives $ps=pt=0$ by directness.\\
Thus $v\in S[p]\oplus T[p]$, and the reverse inclusion is immediate.\\
\hlt{Structural Commentary}: Killing by $p$ turns scalar multiplication modulo $p$ into field scalar multiplication.\\
\hlt{Key Insight}: The layer $M[p]$ counts the number of cyclic primary summands.
\end{proof}
\end{lemma}

\begin{theorem} \label{thm:cyclic-decomposition-primary-pid} \hlt{Cyclic Decomposition of a Primary Module}\\
Let $M$ be f.g. $p$-primary torsion with $o(M)\sim p^e$.\\
Then $M=\bigoplus\limits_{i=1}^m\langle\!\langle v_i\rangle\!\rangle_R$ with $\operatorname{Ann}_R(v_i)=\langle p^{e_i}\rangle$ and $e=e_1\geq e_2\geq\cdots\geq e_m\geq1$.\\
The annihilator chain $\langle p^{e_1}\rangle\subseteq\cdots\subseteq\langle p^{e_m}\rangle$ is intrinsic.
\begin{proof}
\begin{enumerate}[label=(\roman*)]
\setlength{\itemsep}{0pt}
\item \hlt{Existence}: Choose a finite presentation of $M$.\\
Corollary \ref{cor:smith-normal-form-presents-pid-modules} gives $M\cong R^s\oplus\bigoplus\limits_{\ell=1}^tR/\langle d_{\ell}\rangle$, with $d_{\ell}$ nonunit. Since $M$ is torsion, $s=0$.\\
Since $M$ is $p$-primary, $d_{\ell}\sim p^{e_{\ell}}$ with $e_{\ell}\geq1$. Reorder so $e_1\geq\cdots\geq e_m$.\\
Proposition \ref{prop:orders-direct-sum-lcm} gives $o(M)\sim p^{e_1}$, hence $e_1=e$.
\item \hlt{Uniqueness}: Put $\kappa=R/\langle p\rangle$. By Lemma
\ref{lem:p-torsion-layer-vector-space}, each summand $R/\langle p^{e_i}\rangle$ contributes one $\kappa$-dimension to $M[p]$, so $\dim_{\kappa}M[p]=m$. Passing from $M$ to $pM$ replaces each $e_i>1$ by $e_i-1$ and removes each $e_i=1$.\\
Induction on $e$ recovers the exponent multiset, hence the annihilator chain.
\end{enumerate}
\noindent
\hlt{Structural Commentary}: A primary module splits into cyclic strings whose lengths are the exponents $e_i$.\\
\hlt{Key Insight}: The dimensions of the layers $M[p],pM[p],p^2M[p],\ldots$ determine the string lengths.
\end{proof}
\end{theorem}

\begin{proposition} \label{prop:p-primary-layer-counts-recover-exponents} \hlt{Layer Counts Recover Primary Exponents}\\
Let $\kappa=R/\langle p\rangle$ and $\lambda_k(M)\defeq\dim_{\kappa}(p^{k-1}M/p^kM)$ for $k\geq1$.\\
If $M\cong\bigoplus\limits_{i=1}^mR/\langle p^{e_i}\rangle$ with $e_1\geq\cdots\geq e_m\geq1$, then $\lambda_k(M)=\#\{i:e_i\geq k\}$.\\
Hence the number of summands of exact exponent $k$ is $\lambda_k(M)-\lambda_{k+1}(M)$.
\begin{proof} For $C=R/\langle p^e\rangle$, $p^{k-1}C/p^kC$ is $0$ if $k>e$ and is $\kappa$ if $k\leq e$, generated by the class of $p^{k-1}$.\\
Since direct sums commute with $p^{k-1}M$ and $p^kM$, dimensions add, giving $\lambda_k(M)=\#\{i:e_i\geq k\}$.\\
Thus $\lambda_k-\lambda_{k+1}$ counts exactly the summands ending at exponent $k$.\\
\hlt{Structural Commentary}: Successive $p$-layers determine the exponent list.\\
\hlt{Key Insight}: A cyclic summand of order $p^e$ contributes one dimension to exactly the first $e$ layers.
\end{proof}
\end{proposition}

\subsection{Elementary Divisors and Invariant Factors}

\begin{definition} \label{def:elementary-divisors-pid-module} \hlt{Elementary Divisors}\\
For finitely generated torsion $M$, if its primary cyclic summands have annihilators
$\langle p_i^{e_{i,j}}\rangle$,\\
set $\operatorname{ElemDiv}(M)\defeq\{p_i^{e_{i,j}}\}_{i,j}$ as a multiset, up to associates.
\end{definition}

\begin{theorem} \label{thm:primary-cyclic-decomposition-pid} \hlt{Primary Cyclic Decomposition}\\
Every f.g. torsion $R$-module $M$ admits $M=\bigoplus\limits_i\bigoplus\limits_{j=1}^{k_i}\langle\!\langle v_{i,j}\rangle\!\rangle_R$, where $\operatorname{Ann}_R(v_{i,j})=\langle p_i^{e_{i,j}}\rangle$, $p_i$ are distinct nonassociate primes, and $e_{i,1}\geq\cdots\geq e_{i,k_i}\geq1$. For f.g. torsion $M,N$, $M\cong N\iff\operatorname{ElemDiv}(M)=\operatorname{ElemDiv}(N)$.
\begin{proof}
\begin{enumerate}[label=\roman*.]
\setlength{\itemsep}{0pt}
\item[$(\Rightarrow)$] Theorem \ref{thm:primary-decomposition-pid-modules} gives $M=\bigoplus\limits_iM_{p_i}$ with each $M_{p_i}$ f.g. $p_i$-primary. Theorem \ref{thm:cyclic-decomposition-primary-pid} gives $M_{p_i}=\bigoplus\limits_{j=1}^{k_i}\langle\!\langle v_{i,j}\rangle\!\rangle_R$ with $\operatorname{Ann}_R(v_{i,j})=\langle p_i^{e_{i,j}}\rangle$. If $M\cong N$, the same uniqueness statements identify the primary pieces and their exponent chains; hence $\operatorname{ElemDiv}(M)=\operatorname{ElemDiv}(N)$.
\item[$(\Leftarrow)$] If $\operatorname{ElemDiv}(M)=\operatorname{ElemDiv}(N)$, then both modules are isomorphic to the same direct sum $\bigoplus\limits_{q\in\operatorname{ElemDiv}(M)}R/\langle q\rangle$, since a cyclic summand with annihilator $\langle q\rangle$ is $R/\langle q\rangle$ by Theorem
\ref{thm:cyclic-modules-over-pid}. Thus $M\cong N$.
\end{enumerate}
\noindent
\hlt{Structural Commentary}: A torsion module is a direct sum of prime-power cyclic blocks.\\
\hlt{Key Insight}: Elementary divisors are the atoms of the PID module classification.
\end{proof}
\end{theorem}

\begin{proposition} \label{prop:elementary-divisors-direct-sums} \hlt{Elementary Divisors of a Direct Sum}\\
For f.g. torsion $A,B$, $\operatorname{ElemDiv}(A\oplus B)= \operatorname{ElemDiv}(A)\uplus\operatorname{ElemDiv}(B)$, where $\uplus$ is multiset union.
\begin{proof}
If $A=\bigoplus_i A_i$ and $B=\bigoplus_jB_j$ are primary cyclic decompositions, then $A\oplus B=(\bigoplus_i A_i)\oplus(\bigoplus_jB_j)$ is one for $A\oplus B$. Thus the prime-power annihilator multiset is the multiset union.\\
\hlt{Structural Commentary}: Direct sums concatenate cyclic primary blocks.\\
\hlt{Key Insight}: Elementary divisors count blocks with multiplicity.
\end{proof}
\end{proposition}

\begin{definition} \label{def:invariant-factors-pid-module} \hlt{Invariant Factors}\\
For f.g. torsion $M$, write $\operatorname{ElemDiv}(M)=\{p_i^{e_{i,j}}:1\leq j\leq k_i\}_i$ with $e_{i,1}\geq\cdots\geq e_{i,k_i}\geq1$.\\
Let $m=\max_i k_i$, set $e_{i,j}=0$ for $j>k_i$, and define $d_j\defeq\prod_i p_i^{e_{i,j}}$.\\
The nonunit $d_j$, up to associates, are the \hlt{invariant factors}; they satisfy $d_m\mid\cdots\mid d_1$.
\end{definition}

\begin{method} \label{meth:elementary-divisors-to-invariant-factors} \hlt{From Elementary Divisors to Invariant Factors}
\begin{enumerate}[label=(\roman*)]
\setlength{\itemsep}{0pt}
\item Group $\operatorname{ElemDiv}(M)$ by prime: $\{p_i^{e_{i,j}}\}_j$.
\item Reorder each list so $e_{i,1}\geq e_{i,2}\geq\cdots$ and pad by
$e_{i,j}=0$ to common length.
\item Form $d_j=\prod_i p_i^{e_{i,j}}$.
\item Discard unit $d_j$; the rest are the invariant factors from Definition
\ref{def:invariant-factors-pid-module}.
\end{enumerate}
\end{method}

\begin{example} \label{ex:elementary-divisors-to-invariant-factors-polynomial-module} \hlt{Packing Elementary Divisors into Invariant Factors}\\
Over $R=\mathbb{F}[z]$, let $\operatorname{ElemDiv}(M)=\{(z-1)^3,z-1,(z^2+1)^2,z^2+1\}$. The exponent rows are $(3,1)$ for $z-1$ and $(2,1)$ for $z^2+1$, so $d_1=(z-1)^3(z^2+1)^2$ and $d_2=(z-1)(z^2+1)$, with $d_2\mid d_1$. Hence $M\cong R/\langle d_1\rangle\oplus R/\langle d_2\rangle$.\\
\hlt{Structural Commentary}: Elementary divisors separate primes; invariant factors pack rows.\\
\hlt{Key Insight}: Line up equal-height primary strings and multiply across the row.
\end{example}

\begin{theorem} \label{thm:pid-invariant-factor-decomposition} \hlt{Invariant Factor Decomposition}\\
For f.g. torsion $M$, $M\cong\bigoplus_{j=1}^mR/\langle d_j\rangle$ with $d_j\notin R^\times$ and $d_m\mid\cdots\mid d_1$.\\
Writing $\operatorname{Inv}(M)\defeq(d_m,\ldots,d_1)$ up to associates, for f.g. torsion $M,N$, $M\cong N\iff\operatorname{Inv}(M)=\operatorname{Inv}(N)$.
\begin{proof} By Theorem \ref{thm:primary-cyclic-decomposition-pid}, $M\cong\bigoplus_{q\in\operatorname{ElemDiv}(M)}R/\langle q\rangle$.\\
Let $\operatorname{ElemDiv}(M)=\{p_i^{e_{i,j}}:1\leq j\leq k_i\}_i$ with $e_{i,1}\geq\cdots\geq e_{i,k_i}$, pad by $e_{i,j}=0$, and set $d_j=\prod_ip_i^{e_{i,j}}$.\\
For fixed $j$, the nonunit $p_i^{e_{i,j}}$ are pairwise coprime, so Theorem \ref{thm:cyclic-composition-coprime-orders} gives $\bigoplus_{e_{i,j}>0}R/\langle p_i^{e_{i,j}}\rangle\cong R/\langle d_j\rangle$.\\
Hence $M\cong\bigoplus_jR/\langle d_j\rangle$ with $d_m\mid\cdots\mid d_1$.
\begin{enumerate}[label=\roman*.]
\setlength{\itemsep}{0pt}
\item[$(\Rightarrow)$] If $M\cong N$, then $\operatorname{ElemDiv}(M)=\operatorname{ElemDiv}(N)$ by Theorem \ref{thm:primary-cyclic-decomposition-pid}. \\
Grouping by height gives $\operatorname{Inv}(M)=\operatorname{Inv}(N)$.
\item[$(\Leftarrow)$] If $\operatorname{Inv}(M)=\operatorname{Inv}(N)$, by factorization of each $d_j$, $\operatorname{ElemDiv}(M)=\operatorname{ElemDiv}(N)$; QED (Theorem \ref{thm:primary-cyclic-decomposition-pid}).
\end{enumerate}
\hlt{Structural Commentary}: Invariant factors pack prime-power blocks.\\
\hlt{Key Insight}: Same data, different packaging.
\end{proof}
\end{theorem}

\begin{corollary} \label{cor:general-finitely-generated-pid-module-classification} \hlt{Classification of Finitely Generated PID Modules}\\
For f.g. $M$, $M\cong R^r\oplus\bigoplus_{j=1}^mR/\langle d_j\rangle$, where $r\geq0$, $d_j\notin R^\times$, and $d_m\mid\cdots\mid d_1$.\\
The pair $(r,\operatorname{Inv}(M_{\operatorname{tor}}))$, equivalently $(r,\operatorname{ElemDiv}(M_{\operatorname{tor}}))$, is a complete invariant.
\begin{proof} Let $T\defeq M_{\operatorname{tor}}$. Theorem \ref{thm:finitely-generated-pid-free-plus-torsion} gives $M=T\oplus F$ with $F\cong R^r$.\\
Theorem \ref{thm:pid-invariant-factor-decomposition} gives $T\cong\bigoplus_{j=1}^mR/\langle d_j\rangle$ with $d_m\mid\cdots\mid d_1$. Since $T$ is intrinsic and $F\cong M/T$, uniqueness and completeness follow from rank uniqueness and the invariant-factor theorem.\\
\hlt{Structural Commentary}: F.g. PID modules are free plus cyclic torsion.\\
\hlt{Key Insight}: The free rank records the zero-annihilator part; invariant factors record the torsion part.
\end{proof}
\end{corollary}

\begin{perspective} \label{per:invariant-factors-order-and-characteristic-polynomial} \hlt{Largest Invariant Factor and Product}\\
If $M\cong\bigoplus\limits_{j=1}^mR/\langle d_j\rangle$ with $d_m\mid\cdots\mid d_1$, then $\operatorname{Ann}_R(M)=\langle d_1\rangle$ and $\prod\limits_{j=1}^md_j\sim\prod_{q\in\operatorname{ElemDiv}(M)}q$.\\
For $R=\mathbb{F}[z]$ and $M=(\mathbb{F}^n)_A$, after monic normalisation, $d_1=m_A$ and $\prod\limits_jd_j=\chi_A$.
\end{perspective}

\begin{theorem} \label{thm:cyclic-torsion-module-characterization} \hlt{Characterizing Cyclic Torsion Modules}\\
Let $M$ be f.g. torsion with order $\mu\sim\prod_{i=1}^np_i^{e_i}$, $p_i$ distinct nonassociate primes. These are equivalent:
\begin{enumerate}[label=(\roman*)]
\setlength{\itemsep}{0pt}
\item $M$ is cyclic.
\item $\forall i$, the primary component $M_{p_i}$ is cyclic.
\item $\operatorname{ElemDiv}(M)=\{p_i^{e_i}:1\leq i\leq n\}$.
\item $M\cong R/\langle\mu\rangle$.
\end{enumerate}
\begin{proof}
\begin{enumerate}[label=(\roman*).]
\setlength{\itemsep}{0pt}
\item[$(i)$]$\Rightarrow (ii)$ If $M=\langle\!\langle x\rangle\!\rangle_R$, Theorem \ref{thm:cyclic-composition-coprime-orders} gives $M=\bigoplus_i\langle\!\langle x_i\rangle\!\rangle_R$ with $o(x_i)\sim p_i^{e_i}$; hence $M_{p_i}=\langle\!\langle x_i\rangle\!\rangle_R$.
\item[$(ii)$]$\Rightarrow (iii)$ By Theorem
\ref{thm:primary-cyclic-decomposition-pid}, cyclic $M_{p_i}$ of order $p_i^{e_i}$ has sole elementary divisor $p_i^{e_i}$; use Proposition \ref{prop:elementary-divisors-direct-sums}.
\item[$(iii)$]$\Rightarrow (iv)$ Method \ref{meth:elementary-divisors-to-invariant-factors} packs the displayed elementary divisors into one invariant factor $d_1\sim\prod_ip_i^{e_i}\sim\mu$.
\item[$(iv)$]$\Rightarrow (i)$ Since
$R/\langle\mu\rangle=\langle\!\langle 1+\langle\mu\rangle\rangle\!\rangle_R$,
any isomorphism $R/\langle\mu\rangle\xrightarrow{\sim}M$ sends
$1+\langle\mu\rangle$ to a generator of $M$.
\end{enumerate}
\hlt{Structural Commentary}: Cyclic torsion means one string per prime.\\
\hlt{Key Insight}: Multiple elementary divisors for the same prime are the obstruction to one generator.
\end{proof}
\end{theorem}

\begin{definition} \label{def:indecomposable-module} \hlt{Indecomposable Module}\\
A nonzero $R$-module $M$ is \hlt{indecomposable} if $M=A\oplus B\Longrightarrow A=0$ or $B=0$.
\end{definition}

\begin{theorem} \label{thm:indecomposable-fg-torsion-pid-modules} \hlt{Indecomposable Finitely Generated Torsion Modules}\\
Let $0\neq M$ be f.g. torsion. The following are equivalent:
\begin{enumerate}[label=(\roman*)]
\setlength{\itemsep}{0pt}
\item $M$ is indecomposable.
\item $M$ is primary cyclic.
\item $\operatorname{ElemDiv}(M)=\{p^e\}$ for some prime $p$ and $e\geq1$.
\end{enumerate}
\begin{proof}
\begin{enumerate}[label=(\roman*).]
\setlength{\itemsep}{0pt}
\item[$(i)$]$\Rightarrow (ii)$ By Theorem
\ref{thm:primary-cyclic-decomposition-pid}, write $M\cong\bigoplus_{\alpha\in A}C_\alpha$, where $0\neq C_\alpha\cong R/\langle p_\alpha^{e_\alpha}\rangle$.\\
If $\abs{A}\geq2$, then $M\cong C_{\alpha_0}\oplus\bigoplus_{\alpha\neq\alpha_0}C_\alpha$ is a nontrivial decomposition, contradiction.\\
Since $M\neq0$, $\abs{A}=1$; hence $M$ is primary cyclic.
\item[$(ii)$]$\Rightarrow (iii)$ If $M\cong R/\langle p^e\rangle$, then $\operatorname{ElemDiv}(M)=\{p^e\}$.
\item[$(iii)$]$\Rightarrow (i)$ If $M=A\oplus B$, Proposition \ref{prop:elementary-divisors-direct-sums} gives $\{p^e\}=\operatorname{ElemDiv}(A)\sqcup\operatorname{ElemDiv}(B)$.\\
Hence one multiset is empty, so one summand is $0$.
\end{enumerate}
\hlt{Structural Commentary}: Indecomposable torsion modules are precisely the single primary cyclic blocks.\\
\hlt{Key Insight}: Elementary divisors are the indecomposable pieces.
\end{proof}
\end{theorem}

\begin{proposition} \label{prop:prime-divisor-gives-prime-order-submodule} \hlt{Prime Divisors Give Prime-Order Cyclic Submodules}\\
Let $M$ be f.g. torsion with order $\mu$. If $p\mid\mu$ is prime, then $\exists w\in M$ with $o(w)\sim p$; equivalently $Rw\cong R/\langle p\rangle$.
\begin{proof} Write $\mu=pq$. Since $q\notin\operatorname{Ann}_R(M)=\langle\mu\rangle$, choose $v\in M$ with $w\defeq qv\neq0$.\\
Then $pw=\mu v=0$, so $o(w)\mid p$. Since $w\neq0$, $o(w)$ is nonunit; as $p$ is prime, $o(w)\sim p$.\\
\hlt{Structural Commentary}: Every prime in the global order appears in some one-dimensional primary layer.\\
\hlt{Key Insight}: Divide the order by one prime and apply the remaining scalar to an element it does not kill.
\end{proof}
\end{proposition}

